\subsection{Buenas Prácticas de Seguridad}

Para mantener el entorno seguro con medidas concretas y breves:

\begin{itemize}
    \item \textbf{Antimalware en equipos conectados:} motor actualizado, análisis programados.
    \item \textbf{Cifrado en tránsito:} usar TLS 1.2 basado en AES-256; en WiFi, WPA3. 
    \item \textbf{Cifrado en reposo:} proteger respaldos y dispositivos críticos con AES-256; gestionar llaves en un almacén seguro con rotación periódica.
    \item \textbf{VPN institucional:} todo acceso remoto debe ir por la VPN. Aplicar mínimo privilegio y dar de baja credenciales inactivas.
    \item \textbf{Acceso a cuartos de racks:} control físico con cerraduras electrónicas y autenticación biométrica preferente, registro de entradas y política de visitantes.
    \item \textbf{Control de acceso lógico:} mínimo privilegio, cuentas nominales y una robusta auditoría periódica.
    \item \textbf{IDS/IPS:} desplegar en segmentos críticos, con firmas actualizadas y revisión de alertas.
\end{itemize}
